% !TEX root = ../main.tex

\chapter{Additional Materials}

\section{Ontology Definitions}

\todo{Add node definitions, correction edges, and clinical grounding tables if they are too large for the main text.}

\section{Classifier Prompts and Label Mapping}
\label{app:classifier-prompts}

This appendix documents the model-facing labels and prompt structure used by the ontology classifier described in \Cref{sec:safety-layer-classification}. The evaluated default classifier used the prompts below without the optional context-aware few-shot examples implemented for later ablations.

\subsection{Code-to-Node Mapping}

\begin{longtable}{>{\ttfamily}p{0.28\linewidth}>{\ttfamily}p{0.58\linewidth}}
  \caption{Classifier codes and corresponding ontology nodes. Category meanings are given in the prompt excerpt below.}
  \label{tab:classifier-code-map}\\
  \toprule
  Code & Ontology node \\
  \midrule
  \endfirsthead
  \toprule
  Code & Ontology node \\
  \midrule
  \endhead
  HIDE & AVOIDANCE \\
  ESCAPE & INTERRUPTION \\
  DESTROY & VIOLENT\_REVENGE \\
  DENY & SUPPRESSION \\
  REPLAY & TRAUMA\_REPLAY \\
  CONFRONT & BEHAVIORAL\_MASTERY \\
  HELP & SOCIAL\_MASTERY \\
  CHANGE & ENVIRONMENTAL\_MASTERY \\
  RELAX & EMOTIONAL\_MASTERY \\
  POWER & MYTHICAL\_MASTERY \\
  LOOK & NARRATIVE\_SETTING \\
  FEEL & AFFECT\_EXPRESSION \\
  \bottomrule
\end{longtable}

\subsection{System 1 Direct Classifier Prompt}

System 1 asks the model to return exactly one code. The production prompt contains the following category definitions, boundary rules, examples, and response constraint.

\begin{quote}
\small
You are a Clinical Coder for Imagery Rehearsal Therapy (IRT). Classify dream actions into one of 12 codes.

\textbf{Maladaptive codes:}
\texttt{HIDE} for avoidance, hiding, freezing by choice, looking away, passive observation, or choosing to do nothing;
\texttt{ESCAPE} for forcing awake, waking up, abrupt dream ending, or refusing to continue;
\texttt{DESTROY} for brutal excessive violence, torture, mutilation, killing everyone, or uncontrolled rage;
\texttt{DENY} for emotional suppression, claiming to feel nothing, denying distress in a distressing context, or emotional numbing;
\texttt{REPLAY} for re-experiencing trauma without change, rehearsing the nightmare identically, refusing to alter the narrative, or reliving events exactly as they happened.

\textbf{Adaptive codes:}
\texttt{CONFRONT} for fighting back proportionately, running purposefully, taking control, or direct action;
\texttt{HELP} for calling for help, allies appearing, working with others, or asking for assistance;
\texttt{CHANGE} for finding weapons or tools, using objects, favorable environmental change, or manipulating the setting;
\texttt{RELAX} for calming down, deep breathing, accepting peacefully, transforming fear, or staying present;
\texttt{POWER} for magical powers, divine intervention, superpowers, or impossible abilities.

\textbf{Neutral codes:}
\texttt{LOOK} for describing environment, characters, or setting without action;
\texttt{FEEL} for stating feelings without a behavioral strategy.

\textbf{Critical rules:}
controlled defense maps to \texttt{CONFRONT}, while excessive or rageful violence maps to \texttt{DESTROY};
expressing emotion maps to \texttt{FEEL}, while actively disengaging from emotion maps to \texttt{HIDE};
seeking help maps to \texttt{HELP}, using objects or changing the setting maps to \texttt{CHANGE}, and direct physical action maps to \texttt{CONFRONT};
negated avoidance, such as deciding not to run, maps to \texttt{CONFRONT};
imagined shields or protection map to \texttt{POWER} or \texttt{CHANGE}, breathing or self-reassurance maps to \texttt{RELAX}, and trying not to feel maps to \texttt{DENY};
engaged emotional regulation maps to \texttt{RELAX}, while emotional shutdown in a distressing context maps to \texttt{DENY};
scene description maps to \texttt{LOOK}, while reliving trauma identically maps to \texttt{REPLAY}.

\textbf{Examples:}
``I hide under the bed'' $\rightarrow$ \texttt{HIDE};
``I wake myself up'' $\rightarrow$ \texttt{ESCAPE};
``I kill everyone brutally'' $\rightarrow$ \texttt{DESTROY};
``I don't feel anything about it'' $\rightarrow$ \texttt{DENY};
``It all happens exactly the same way again'' $\rightarrow$ \texttt{REPLAY};
``I punch the attacker'' $\rightarrow$ \texttt{CONFRONT};
``I call for help'' $\rightarrow$ \texttt{HELP};
``I find a sword'' $\rightarrow$ \texttt{CHANGE};
``I take a deep breath'' $\rightarrow$ \texttt{RELAX};
``I use magic to fly'' $\rightarrow$ \texttt{POWER};
``There was a dark forest'' $\rightarrow$ \texttt{LOOK};
``I felt terrified'' $\rightarrow$ \texttt{FEEL}.

\textbf{Response format:}
Respond with only one code from \texttt{HIDE}, \texttt{ESCAPE}, \texttt{DESTROY}, \texttt{DENY}, \texttt{REPLAY}, \texttt{CONFRONT}, \texttt{HELP}, \texttt{CHANGE}, \texttt{RELAX}, \texttt{POWER}, \texttt{LOOK}, or \texttt{FEEL}. No explanations.
\end{quote}

\subsection{System 2 Structured Reasoning Prompt}

System 2 uses the same code definitions, rules, and examples as System 1. The only intended change is the response format: the model is instructed to return a JSON object with the fields \texttt{context\_summary}, \texttt{clinical\_rationale}, and \texttt{classification}. The \texttt{classification} field must contain one of the same 12 codes listed in \Cref{tab:classifier-code-map}.

\section{Example Safety Constraint Message}
\label{app:example-constraint-message}

The following abbreviated example illustrates the kind of constraint message injected into the response-generator prompt when the critic detects a maladaptive segment alongside adaptive material. It corresponds to the mixed example in \Cref{fig:worked-end-to-end-example}. Line breaks and wording are preserved in simplified form for readability.

\begin{quote}
\small
\texttt{[SAFETY SYSTEM - INTERVENTION REQUIRED]}

\texttt{VALIDATE:} The user showed positive elements: ``I call my brother for help''. Acknowledge these first.

\texttt{CORRECT:} User demonstrated \texttt{VIOLENT\_REVENGE}; this is a \texttt{CRITICAL} risk.

Risk description: uncontrolled, chaotic violence that destroys the scene rather than resolving it.

Clinical guideline: guide toward controlled, purposeful action instead of destructive chaos.

\texttt{INSTRUCTION:} Gently redirect using this approach: What focused action could resolve this without destroying everything? How could the user stop the threat without excessive violence?

Alternative therapeutic directions:
\texttt{BEHAVIORAL\_MASTERY}: controlled, purposeful action;
\texttt{SOCIAL\_MASTERY}: seeking or working with support.

\texttt{IMPORTANT:} Do not validate, restate, or encourage the maladaptive response, here ``I burn the whole city down''. Use supportive language, focus the correction on this issue, and do not shame or criticize the user's suggestion.
\end{quote}

\section{Benchmark Details}

\todo{Add dataset-generation prompts, scenario distributions, and extended result tables.}

\section{Expert Annotation Materials}

\todo{Add questionnaire, rating instructions, or annotation protocol excerpts as appropriate.}
